% ============================================================
\chapter{Int\'egrale Stochastique d'It\^o}
\label{ch:ito_integrale}
% ============================================================

\section{Motivation}

En 1942, \`a Kyoto, le math\'ematicien Kiyosi It\^o s'attaque \`a un probl\`eme fondamental~: comment d\'efinir une int\'egrale par rapport au mouvement brownien~? Les trajectoires du brownien sont continues mais nulle part diff\'erentiables, et leur variation totale est infinie sur tout intervalle. L'int\'egrale de Stieltjes classique ne s'applique donc pas. It\^o contourne l'obstacle par une construction ing\'enieuse~: il approche l'int\'egrande par des processus \'etag\'es et montre que la limite existe dans $L^2$, gr\^ace \`a l'isom\'etrie d'It\^o. Le choix du point d'\'evaluation \`a \emph{gauche} de chaque intervalle (contrairement \`a Stratonovich, qui prend le milieu) donne \`a l'int\'egrale d'It\^o la propri\'et\'e d'\^etre une martingale --- propri\'et\'e cruciale pour la finance et les probabilit\'es.

L'objectif de ce chapitre est de donner un sens rigoureux \`a l'int\'egrale
\[
  \int_0^t H_s \, dB_s,
\]
o\`u $(B_t)$ est un mouvement brownien et $(H_t)$ un processus adapt\'e. Comme
les trajectoires du brownien sont \`a variation totale infinie
(th\'eor\`eme~\ref{thm:var_totale}), cette int\'egrale ne peut pas \^etre d\'efinie au
sens de Stieltjes. La construction d'It\^o proc\`ede par approximation.

\section{Construction pour les processus \'etag\'es}

\begin{definition}[Processus \'etag\'e]
Un processus $(H_t)_{t \in [0,T]}$ est \emph{\'etag\'e} (ou \emph{simple}, ou
\emph{en escalier}) s'il existe une partition $0 = t_0 < t_1 < \cdots < t_n = T$
et des v.a.\ born\'ees $\mathcal{F}_{t_k}$-mesurables $H_k$ telles que~:
\[
  H_t = H_0 \mathbf{1}_{\{0\}}(t) + \sum_{k=0}^{n-1} H_k \mathbf{1}_{(t_k, t_{k+1}]}(t).
\]
On note $\mathcal{S}$ l'espace des processus \'etag\'es.
\end{definition}

\begin{definition}[Int\'egrale d'It\^o pour les processus \'etag\'es]
Pour $H \in \mathcal{S}$, on d\'efinit~:
\[
  I(H) = \int_0^T H_s \, dB_s := \sum_{k=0}^{n-1} H_k (B_{t_{k+1}} - B_{t_k}).
\]
\end{definition}

\begin{attention}[title=\textbf{\'Evaluation \`a gauche}]
Dans la d\'efinition d'It\^o, l'int\'egrande $H_k$ est \'evalu\'e au temps
\emph{gauche} $t_k$ de l'intervalle $(t_k, t_{k+1}]$. Ce choix est crucial
pour la propri\'et\'e de martingale. Un choix diff\'erent (par exemple au milieu,
comme dans l'int\'egrale de Stratonovich) donne un r\'esultat diff\'erent.
\end{attention}

\begin{proposition}[Propri\'et\'es de l'int\'egrale pour les processus \'etag\'es]
\label{prop:prop_etage}
Pour $H, G \in \mathcal{S}$ et $a, b \in \R$~:
\begin{enumerate}
  \item \textbf{Lin\'earit\'e}~: $I(aH + bG) = aI(H) + bI(G)$.
  \item \textbf{Esp\'erance nulle}~: $\E[I(H)] = 0$.
  \item \textbf{Isom\'etrie d'It\^o}~:
    $\E[I(H)^2] = \E\Bigl[\int_0^T H_s^2 \, ds\Bigr]$.
  \item \textbf{Propri\'et\'e de martingale}~: le processus
    $t \mapsto \int_0^t H_s \, dB_s$ est une martingale.
\end{enumerate}
\end{proposition}

\begin{proof}[Preuve de l'isom\'etrie d'It\^o]
\begin{align*}
  \E[I(H)^2] &= \E\Bigl[\Bigl(\sum_{k=0}^{n-1} H_k \Delta B_k\Bigr)^2\Bigr] \\
  &= \sum_{k=0}^{n-1} \E[H_k^2 (\Delta B_k)^2]
    + 2 \sum_{j < k} \E[H_j \Delta B_j H_k \Delta B_k].
\end{align*}
Les termes crois\'es s'annulent car, pour $j < k$~:
\[
  \E[H_j \Delta B_j H_k \Delta B_k]
  = \E[\underbrace{H_j \Delta B_j H_k}_{\mathcal{F}_{t_k}\text{-mesurable}}
  \underbrace{\E[\Delta B_k \mid \mathcal{F}_{t_k}]}_{=0}] = 0.
\]
Pour les termes diagonaux~:
\[
  \E[H_k^2 (\Delta B_k)^2] = \E[H_k^2 \E[(\Delta B_k)^2 \mid \mathcal{F}_{t_k}]]
  = \E[H_k^2 (t_{k+1} - t_k)].
\]
Donc $\E[I(H)^2] = \sum_k \E[H_k^2 (t_{k+1} - t_k)] = \E[\int_0^T H_s^2 \, ds]$.
\end{proof}

\section{Extension aux processus $L^2$}

\begin{definition}[Espace $\mathcal{H}^2$]
On d\'efinit l'espace des int\'egrandes admissibles~:
\[
  \mathcal{H}^2 = \mathcal{H}^2([0, T]) = \Bigl\{H \text{ adapt\'e}
  : \E\Bigl[\int_0^T H_s^2 \, ds\Bigr] < \infty \Bigr\}.
\]
Muni de la norme $\norm{H}_{\mathcal{H}^2} = \bigl(\E[\int_0^T H_s^2 \, ds]\bigr)^{1/2}$,
c'est un espace de Hilbert.
\end{definition}

\begin{theoreme}[Extension de l'int\'egrale d'It\^o]
\label{thm:extension_ito}
L'application lin\'eaire isom\'etrique $I : \mathcal{S} \to L^2(\Omega)$ s'\'etend
de mani\`ere unique en une isom\'etrie $I : \mathcal{H}^2 \to L^2(\Omega)$.
L'int\'egrale \'etendue, not\'ee
\[
  \int_0^T H_s \, dB_s = I(H),
\]
v\'erifie les m\^emes propri\'et\'es que pour les processus \'etag\'es~: lin\'earit\'e,
esp\'erance nulle, isom\'etrie d'It\^o.
\end{theoreme}

\begin{proof}[Id\'ee de la preuve]
L'espace $\mathcal{S}$ est dense dans $\mathcal{H}^2$ (on peut approcher tout
processus adapt\'e de carr\'e int\'egrable par des processus \'etag\'es). L'isom\'etrie
d'It\^o garantit que $I$ est une isom\'etrie de $(\mathcal{S}, \norm{\cdot}_{\mathcal{H}^2})$
dans $(L^2(\Omega), \norm{\cdot}_{L^2})$. Par compl\'etude de $L^2(\Omega)$,
$I$ s'\'etend de mani\`ere unique.
\end{proof}

\section{Propri\'et\'es de l'int\'egrale d'It\^o}

\begin{theoreme}[Propri\'et\'es fondamentales]
\label{thm:prop_ito}
Soit $H \in \mathcal{H}^2([0,T])$. Le processus $M_t = \int_0^t H_s \, dB_s$
v\'erifie~:
\begin{enumerate}
  \item \textbf{Continuit\'e}~: $(M_t)$ a une version \`a trajectoires continues.
  \item \textbf{Martingale}~: $(M_t)_{0 \leq t \leq T}$ est une martingale de carr\'e
    int\'egrable par rapport \`a $(\mathcal{F}_t)$.
  \item \textbf{Isom\'etrie d'It\^o}~: $\E[M_t^2] = \E[\int_0^t H_s^2 \, ds]$.
  \item \textbf{Variation quadratique}~: $\langle M \rangle_t = \int_0^t H_s^2 \, ds$.
\end{enumerate}
\end{theoreme}

\begin{formules}[title=\textbf{Propri\'et\'es cl\'es de l'int\'egrale d'It\^o}]
Pour $M_t = \int_0^t H_s \, dB_s$ avec $H \in \mathcal{H}^2$~:
\begin{align*}
  \E[M_t] &= 0 \\
  \E[M_t^2] &= \E\Bigl[\int_0^t H_s^2 \, ds\Bigr]
    \quad \text{(isom\'etrie d'It\^o)} \\
  \langle M \rangle_t &= \int_0^t H_s^2 \, ds \\
  dM_t &= H_t \, dB_t
\end{align*}
\end{formules}

\section{Extension aux int\'egrandes localement $L^2$}

\begin{definition}[Int\'egrale d'It\^o g\'en\'eralis\'ee]
Pour un processus adapt\'e $H$ avec $\int_0^T H_s^2 \, ds < \infty$ p.s.\ (mais
pas n\'ecessairement $\E[\int_0^T H_s^2 \, ds] < \infty$), on d\'efinit l'int\'egrale
d'It\^o par localisation~: on pose $\tau_n = \inf\{t : \int_0^t H_s^2 \, ds \geq n\}$
et $\int_0^t H_s \, dB_s = \lim_{n \to \infty} \int_0^{t \wedge \tau_n} H_s \, dB_s$.

Le processus r\'esultant est une \emph{martingale locale} continue.
\end{definition}

\section{Int\'egrale de Stratonovich}

\begin{definition}[Int\'egrale de Stratonovich]
L'\emph{int\'egrale de Stratonovich} est d\'efinie par~:
\[
  \int_0^T H_s \circ dB_s = \lim_{\abs{\Pi} \to 0}
  \sum_{k=0}^{n-1} \frac{H_{t_k} + H_{t_{k+1}}}{2} (B_{t_{k+1}} - B_{t_k}).
\]
\end{definition}

\begin{proposition}[Relation It\^o--Stratonovich]
Si $H_t = f(B_t)$ avec $f \in C^1(\R)$, alors~:
\[
  \int_0^T f(B_s) \circ dB_s = \int_0^T f(B_s) \, dB_s
  + \frac{1}{2} \int_0^T f'(B_s) \, ds.
\]
\end{proposition}

\begin{remarque}
L'int\'egrale de Stratonovich ob\'eit aux r\`egles classiques du calcul diff\'erentiel
(pas de terme correctif d'It\^o) mais n'est pas une martingale en g\'en\'eral.
L'int\'egrale d'It\^o est privil\'egi\'ee en finance (propri\'et\'e de martingale),
tandis que l'int\'egrale de Stratonovich est naturelle en physique (invariance
par changement de coordonn\'ees).
\end{remarque}

\section{Exemples de calcul}

\begin{exemple}[$\int_0^t B_s \, dB_s$]
Pour une partition $\Pi = \{0 = t_0 < t_1 < \cdots < t_n = t\}$~:
\begin{align*}
  \sum_{k=0}^{n-1} B_{t_k}(B_{t_{k+1}} - B_{t_k})
  &= \sum_{k=0}^{n-1} \Bigl[\frac{1}{2}(B_{t_{k+1}}^2 - B_{t_k}^2)
    - \frac{1}{2}(B_{t_{k+1}} - B_{t_k})^2\Bigr] \\
  &= \frac{1}{2} B_t^2 - \frac{1}{2} \sum_{k=0}^{n-1} (\Delta B_k)^2.
\end{align*}
En passant \`a la limite et utilisant $\sum (\Delta B_k)^2 \to t$ (variation
quadratique)~:
\[
  \int_0^t B_s \, dB_s = \frac{1}{2} B_t^2 - \frac{1}{2} t.
\]
\end{exemple}

\begin{attention}[title=\textbf{Le terme $-t/2$}]
Le r\'esultat $\int_0^t B_s \, dB_s = \frac{1}{2}B_t^2 - \frac{1}{2}t$ montre que
le calcul stochastique diff\`ere du calcul classique. En calcul ordinaire, on aurait
$\int_0^t x \, dx = x^2/2$. Le terme $-t/2$ est le \og terme correctif d'It\^o \fg,
cons\'equence de la variation quadratique non nulle du brownien.
\end{attention}

\begin{exemple}[$\int_0^t e^{B_s - s/2} \, dB_s$]
Posons $M_t = e^{B_t - t/2}$. On v\'erifiera au chapitre suivant (formule d'It\^o)
que $dM_t = M_t \, dB_t$, donc~:
\[
  \int_0^t e^{B_s - s/2} \, dB_s = M_t - M_0 = e^{B_t - t/2} - 1.
\]
\end{exemple}

\section{Exercices}

\begin{exercice}
Calculer $\E[(\int_0^1 B_s \, dB_s)^2]$ en utilisant l'isom\'etrie d'It\^o.
\end{exercice}

\begin{exercice}
Soit $f \in C^1(\R)$. Montrer que $\int_0^t f(s) \, dB_s$ est une v.a.\ gaussienne
et calculer sa loi.
\end{exercice}

\begin{exercice}
Calculer $\int_0^t s \, dB_s$ et v\'erifier que c'est une v.a.\ gaussienne de
variance $t^3/3$.
\end{exercice}

\begin{exercice}
Montrer que $\int_0^t B_s^2 \, dB_s = \frac{1}{3} B_t^3 - \int_0^t B_s \, ds$
en utilisant la formule d'It\^o (anticiper le chapitre suivant, ou d\'emontrer
directement par approximation).
\end{exercice}

\begin{exercice}
Soit $H_t = \text{sgn}(B_t)$ (signe du brownien). V\'erifier que $H \in \mathcal{H}^2$
et que $\int_0^t \text{sgn}(B_s) \, dB_s = \abs{B_t} - L_t^0$
(formule de Tanaka).
\end{exercice}

\begin{exercice}
Comparer les int\'egrales d'It\^o et de Stratonovich de $\int_0^t B_s \, dB_s$
et v\'erifier la formule de correction.
\end{exercice}
